% Official solution, 2019 DA solutions pp. 1-8; the source's equation
% numbers (1.1)-(1.15) are kept.

\begin{description}
\item[a)] From the light curve plot the Modified Julian Dates can be read
  with an error of about 0.5. The peak of the maximum brightness is very
  narrow and obviously placed at $\mathrm{MJD}_0 = 56\,520.5$ with a value of
  $m_0 = 4.5^m$, therefore:
  \[ \mathrm{MJD}_0 = 56\,520.5 \pm 0.5 \]
  \hfill(2 p)
  \[ m_0 = 4.5^m \pm 0.05^m \]
  \hfill(1 p)

\item[b)] The brightness values of $2^m$ and $3^m$ decline are
  $6.5^m \pm 0.05^m$ and $7.5^m \pm 0.05^m$, respectively. \hfill(2 p)

  The corresponding Modified Julian Dates are $\mathrm{MJD}_2$ and
  $\mathrm{MJD}_3$. Because of the poorly defined slopes on the light curve
  around these dates, their acceptable error is larger than in other parts of
  the light curve, let say it is about $1^d$, so:
  \[ \mathrm{MJD}_2 = 56\,531.5 \pm 1, \qquad
     \mathrm{MJD}_3 = 56\,543.5 \pm 1 \]
  \hfill(2 p)
  \[ t_2 = 11^d \pm 1^d, \qquad t_3 = 23^d \pm 1^d \]
  \hfill(2 p)

\item[c)] The text of this part does not ask for calculating the individual
  errors of the formulae, but it is worth estimating them here, just for the
  sake of completeness. (Students won't do it.)
  \begin{description}
  \item[(a)] The form of the function is
    \[ M = a + b \arctan\frac{c - \log t_2}{d}, \qquad
       a = -7.92,\ b = -0.81,\ c = 1.32,\ d = 0.23,
       \tag*{(1.1)} \]
    so its derivative:
    \[ M' = -\frac{b}{d \log(10)
       \left[ 1 + \left( \frac{c - \log t_2}{d} \right)^2 \right]}
       \frac{1}{t_2}
       \;\rightarrow\;
       \Delta M = -\frac{b}{d \log(10)
       \left[ 1 + \left( \frac{c - \log t_2}{d} \right)^2 \right]}
       \frac{\Delta t_2}{t_2}
       \tag*{(1.2)} \]
    The value and its error calculated from the formulae above (error is not
    necessary):
    \[ M_0^{(a)} = -8.63^m \pm 0.06^m \]
    \hfill(1 p)
  \item[(b)] The form of the function is
    \[ M = a + b \log t_2, \qquad a = -11.32,\ b = 2.55
       \tag*{(1.3)} \]
    so its derivative:
    \[ M' = \frac{b}{t_2 \log(10)}
       \;\rightarrow\;
       \Delta M = \frac{b}{\log(10)} \frac{\Delta t_2}{t_2}
       \tag*{(1.4)} \]
    The value and its error calculated from the formulae above (error is not
    necessary):
    \[ M_0^{(b)} = -8.66^m \pm 0.10^m \]
    \hfill(1 p)
  \item[(c)] The form of the function is
    \[ M = a + b \log t_3, \qquad a = -11.99,\ b = 2.54
       \tag*{(1.5)} \]
    so its derivative:
    \[ M' = \frac{b}{t_3 \log(10)}
       \;\rightarrow\;
       \Delta M = \frac{b}{\log(10)} \frac{\Delta t_3}{t_3}
       \tag*{(1.6)} \]
    The value and its error calculated from the formulae above (error is not
    necessary):
    \[ M_0^{(c)} = -8.53^m \pm 0.05^m \]
    \hfill(1 p)
  \end{description}

  The standard deviation of a dataset can be calculated as
  \[ \sigma = \sqrt{ \frac{\sum\limits_{i=1}^{n} (x_i - \bar{x})^2}{n-1} },
     \tag*{(1.7)} \]
  where $n$ is the number of data points, $x_i$ is the $i$th individual data
  value and $\bar{x}$ is their mean,
  $\bar{x} = (x_1 + x_2 + \ldots + x_n)/n$.

  The mean and the standard deviation of the three absolute maximum
  brightness values:
  \[ M_0 = -8.61^m \pm 0.07^m \]
  \hfill(2 p)

\item[d)] The color excess $E(B-V)$ is the difference between the observed
  color index of the star and the intrinsic color index predicted from its
  spectral type:
  \[ E(B-V) = (B-V) - (B-V)_0 = A_B - A_V
     \tag*{(1.8)} \]
  The total extinction is quantified by $A_V$ (at 5550 \AA). The ratio of
  total-to-selective extinction:
  \[ R = \frac{A_V}{E(B-V)}
     \;\rightarrow\; A_V = R\,E(B-V), \text{ where } R = 3.1
     \tag*{(1.9)} \]
  \hfill(2 p)

  With the given value and error of $E(B-V)$:
  \[ A_V = 3.1 \times E(B-V) = 3.1 \times (0.184^m \pm 0.035^m)
     = 0.57^m \pm 0.11^m \]
  \hfill(2 p)

\item[e)] According to the formula for the distance modulus:
  \[ m_V - M_V = -5 + 5\log d + A_V \;\rightarrow\;
     \tag*{(1.10)} \]
  \hfill(1 p)
  \[ \log d = \frac{m_V - M_V + 5 - A_V}{5} \;\rightarrow\;
     \tag*{(1.11)} \]
  \[ d = 10^{(m_V - M_V + 5 - A_V)/5},
     \tag*{(1.12)} \]
  \hfill(2 p)\\
  where the distance $d$ is in parsecs.

  Since $\Delta a^x / \Delta x = a^x \ln a$, therefore the error of the
  distance $d$:
  \[ \Delta d = 10^{(m_V - M_V + 5 - A_V)/5} \times \ln 10 \times
     \Delta\left( (m_V - M_V + 5 - A_V)/5 \right) \]
  \hfill(2 p)

  The error of $(m_V - M_V + 5 - A_V)/5$ can be estimated with the sum of
  the errors of $m_V$, $M_V$ and $A_V$, so:
  \[ \Delta\left( (m_V - M_V + 5 - A_V)/5 \right) \approx 0.05 \]
  \hfill(2 p)

  With the data:
  \[ d = 3220\ \mathrm{pc} \quad\text{and}\quad \Delta d = 338\ \mathrm{pc}, \]
  \hfill(2 p)\\
  so the distance to the nova:
  \[ d \approx 3.2 \pm 0.3\ \mathrm{kpc} \]
  \hfill(2 p)

\item[f)] The well known Doppler formula between the wavelength displacement
  and radial velocity:
  \[ \frac{\lambda - \lambda_0}{\lambda_0} = \frac{\Delta\lambda}{\lambda_0}
     = \frac{v_r}{c}
     \;\rightarrow\; v_r = \frac{\Delta\lambda}{\lambda_0}\,c,
     \tag*{(1.13)} \]
  \hfill(1 p)\\
  where $\lambda$ is the measured wavelength of the line feature, $\lambda_0$
  is the rest wavelength of the line, $c$ is the speed of light, and $v_r$ is
  the radial velocity to be calculated.

  The wavelength of the P Cygni absorption peak should be extracted from the
  figure with an error of about 1 \AA. The main part of this error is coming
  from the ``definition'' of the peak position of the Gaussian-like profile.
  This will result in inaccuracy of about
  $\Delta v_r = \pm 50\ \mathrm{km\,s^{-1}}$ in radial velocities. \hfill(1 p)

  The wavelengths and radial velocities should be something like these:
  \begin{center}
  \begin{tabular}{ccc}
    \hline
    MJD & WL & RV \\
    \hline
    56518.986 & 6527 & $-1636$ \\
    56519.813 & 6531 & $-1454$ \\
    56520.843 & 6534 & $-1317$ \\
    56521.835 & 6537 & $-1179$ \\
    56522.829 & 6542 & $-951$ \\
    56523.827 & 6544 & $-860$ \\
    \hline
  \end{tabular}
  \end{center}
  \hfill(12 p)

  Radial velocities within the range of $\pm 50\ \mathrm{km\,s^{-1}}$ of the
  RV values listed in the table should be given full marks, but velocities in
  the range of $\pm 100\ \mathrm{km\,s^{-1}}$ are still acceptable with half
  marks.

\item[g)] See the attached figure as an example for the acceptable solution.
  The plotted data are taken from the table above. For the sake of simplicity
  the absolute values of the radial velocities have been used for making the
  graph. \hfill(6 p)
  % FIGURE NEEDED: |v_r| (km/s) of the H-alpha absorption components vs MJD
  % (56519-56524), six points falling along a straight line with fitted
  % dashed line and hatched trapezoid under the curve; annotations "4.841
  % days = 418 262 s" and "1248 km/s"; 2019 DA solutions, page 7.

\item[h)] It is obvious from the plot, that the radial velocities lie along a
  straight line. To estimate the size of the expanding envelope we need to
  calculate the area below the $t - v_r$ graph between the first and last
  date. Hence the graph is a straight line, this is very simple: we have to
  determine the area of the hatched region which is a trapezoid.

  If the two bases and the height of the trapezoid are $a$, $b$, and $m$,
  respectively, then the area of the trapezoid is:
  \[ T = \frac{a + b}{2}\,m
     \tag*{(1.14)} \]
  \hfill(3 p)

  In our case $a = v_{r1}$, $b = v_{r6}$, and $m = t_6 - t_1$. \hfill(1 p)

  We could use the fitted line (dashed) as the upper side of the trapezoid,
  but this would be a bit complicated --- because of the difficulties of the
  fitting process ---, and not necessary at all. Instead of this we use the
  line connecting the first and last radial velocity points as this runs very
  close to the fitted line.
  \[ R = \int_{t_1}^{t_6} |v_r|\,dt
     \approx \frac{(v_{r1} + v_{r6})}{2}\,(t_6 - t_1)
     \approx 3.5\ \mathrm{AU} \]
  The result: $R \approx 3.5$ AU \hfill(3 p)

\item[i)] The apparent angular diameter of the spherical envelope seen from
  the Earth:
  \[ \vartheta = 2 \times \arctan\left( \frac{R}{d} \right)
     \approx 2\,\frac{R}{d}
     \tag*{(1.15)} \]
  \hfill(2 p)

  Using the values of $d \approx 3.2$ kpc and $R \approx 3.5$ AU, 5 days
  after the discovery the angular diameter of the envelope is:
  \[ \vartheta = 0.0022'' = 2.2\ \mathrm{mas} \]
  \hfill(2 p)

  A less formal solution:
  \begin{itemize}
  \item By definition a parsec (1 pc) is the distance from the Sun to an
    astronomical object that has a parallax angle of one arcsecond, i.e.\ it
    represents the distance at which the radius of Earth's orbit (1 AU)
    subtends an angle of one arcsecond.
  \item Because of the very small angles the distance is a linear function of
    the parallax. This means that the radius $R \approx 3.5$ AU of the
    envelope subtends one arcsecond viewing from a distance of
    $d \approx 3.5$ pc, and one milliarcsecond from a distance of
    $1000d \approx 3500\ \mathrm{pc} = 3.5\ \mathrm{kpc}$.
  \item Since this value is close to the distance of Nova Del 2013 determined
    earlier, we can conclude that the apparent angular diameter of the
    spherical shape envelope 5 days after the discovery was about 2
    milliarcseconds.
  \end{itemize}

% FIGURE NEEDED: annotated Johnson V light curve of Nova Del 2013, m_V vs
% MJD (56510-56610), marking MJD_0 = 56520.5 (m_0 = 4.5), MJD_2 = 56531.5
% (m_2 = 6.5), MJD_3 = 56543.5 (m_3 = 7.5), t_2 = 11 d, t_3 = 23 d;
% 2019 DA solutions, page 5.
% FIGURE NEEDED: stacked spectra of Nova Del 2013 around the H-alpha line,
% relative flux + constant vs wavelength 6450-6730 A, six epochs t_0 - 2 d
% ... t_0 + 3 d with P Cygni absorption components; 2019 DA solutions,
% page 6.
\end{description}
